\documentclass[10pt,a4paper]{amsart}
\usepackage{amsmath,amssymb,amsthm,graphicx}
% Retained minimal preamble for the Stacks Project chapter more-morphisms.tex (More on Morphisms).
% Verbatim excerpt of the chapter preamble: Xy-pic setup (native preamble lines 14-18) and the native
% theorem declarations with the mathematical macros (native preamble lines 154-217). Copied mechanically;
% nothing added, removed or redefined. The wrapper class and the amsmath/amssymb/amsthm/graphicx base
% packages are supplied by the assembler, so the class/externaldocument part of the native preamble is omitted.
% For commutative diagrams we use Xy-pic
\usepackage[all]{xy}

% We use 2cell for 2-commutative diagrams.
\xyoption{2cell}
% Theorem environments.
%
\theoremstyle{plain}
\newtheorem{theorem}[subsection]{Theorem}
\newtheorem{proposition}[subsection]{Proposition}
\newtheorem{lemma}[subsection]{Lemma}

\theoremstyle{definition}
\newtheorem{definition}[subsection]{Definition}
\newtheorem{example}[subsection]{Example}
\newtheorem{exercise}[subsection]{Exercise}
\newtheorem{situation}[subsection]{Situation}

\theoremstyle{remark}
\newtheorem{remark}[subsection]{Remark}
\newtheorem{remarks}[subsection]{Remarks}

\numberwithin{equation}{subsection}

% Macros
%
\def\lim{\mathop{\mathrm{lim}}\nolimits}
\def\colim{\mathop{\mathrm{colim}}\nolimits}
\def\Spec{\mathop{\mathrm{Spec}}}
\def\Hom{\mathop{\mathrm{Hom}}\nolimits}
\def\Ext{\mathop{\mathrm{Ext}}\nolimits}
\def\SheafHom{\mathop{\mathcal{H}\!\mathit{om}}\nolimits}
\def\SheafExt{\mathop{\mathcal{E}\!\mathit{xt}}\nolimits}
\def\Sch{\mathit{Sch}}
\def\Mor{\mathop{\mathrm{Mor}}\nolimits}
\def\Ob{\mathop{\mathrm{Ob}}\nolimits}
\def\Sh{\mathop{\mathit{Sh}}\nolimits}
\def\NL{\mathop{N\!L}\nolimits}
\def\CH{\mathop{\mathrm{CH}}\nolimits}
\def\proetale{{pro\text{-}\acute{e}tale}}
\def\etale{{\acute{e}tale}}
\def\QCoh{\mathit{QCoh}}
\def\Ker{\mathop{\mathrm{Ker}}}
\def\Im{\mathop{\mathrm{Im}}}
\def\Coker{\mathop{\mathrm{Coker}}}
\def\Coim{\mathop{\mathrm{Coim}}}

% Boxtimes
%
\DeclareMathSymbol{\boxtimes}{\mathbin}{AMSa}{"02}

%
% Macros for moduli stacks/spaces
%
\def\QCohstack{\mathcal{QC}\!\mathit{oh}}
\def\Cohstack{\mathcal{C}\!\mathit{oh}}
\def\Spacesstack{\mathcal{S}\!\mathit{paces}}
\def\Quotfunctor{\mathrm{Quot}}
\def\Hilbfunctor{\mathrm{Hilb}}
\def\Curvesstack{\mathcal{C}\!\mathit{urves}}
\def\Polarizedstack{\mathcal{P}\!\mathit{olarized}}
\def\Complexesstack{\mathcal{C}\!\mathit{omplexes}}
% \Pic is the operator that assigns to X its picard group, usage \Pic(X)
% \Picardstack_{X/B} denotes the Picard stack of X over B
% \Picardfunctor_{X/B} denotes the Picard functor of X over B
\def\Pic{\mathop{\mathrm{Pic}}\nolimits}
\def\Picardstack{\mathcal{P}\!\mathit{ic}}
\def\Picardfunctor{\mathrm{Pic}}
\def\Deformationcategory{\mathcal{D}\!\mathit{ef}}
\usepackage{cleveref}
\crefname{originallemma}{original lemma}{original lemmas}
\Crefname{originallemma}{Original lemma}{Original lemmas}
\makeatletter
\newlabel{lemma-dominate-fppf}{{05WN}{source}{}{sourcepointer-1}{}}
\newlabel{lemma-dominate-fppf@cref}{{[originallemma][0][]05WN}{[1][0][]source}{}{}{}}
\makeatother
\crefname{originallemma}{original lemma}{original lemmas}
\Crefname{originallemma}{Original lemma}{Original lemmas}
\makeatletter
\newlabel{lemma-finite-morphism-relative-pseudo-coherence}{{09UK}{source}{}{sourcepointer-2}{}}
\newlabel{lemma-finite-morphism-relative-pseudo-coherence@cref}{{[originallemma][0][]09UK}{[1][0][]source}{}{}{}}
\makeatother
\crefname{originallemma}{original lemma}{original lemmas}
\Crefname{originallemma}{Original lemma}{Original lemmas}
\makeatletter
\newlabel{lemma-flat-finite-presentation-pseudo-coherent}{{0695}{source}{}{sourcepointer-3}{}}
\newlabel{lemma-flat-finite-presentation-pseudo-coherent@cref}{{[originallemma][0][]0695}{[1][0][]source}{}{}{}}
\makeatother
\crefname{originallemma}{original lemma}{original lemmas}
\Crefname{originallemma}{Original lemma}{Original lemmas}
\makeatletter
\newlabel{lemma-permanence-pseudo-coherent}{{0683}{source}{}{sourcepointer-4}{}}
\newlabel{lemma-permanence-pseudo-coherent@cref}{{[originallemma][0][]0683}{[1][0][]source}{}{}{}}
\makeatother
\crefname{originallemma}{original lemma}{original lemmas}
\Crefname{originallemma}{Original lemma}{Original lemmas}
\makeatletter
\newlabel{lemma-pull-relative-pseudo-coherent}{{09US}{source}{}{sourcepointer-5}{}}
\newlabel{lemma-pull-relative-pseudo-coherent@cref}{{[originallemma][0][]09US}{[1][0][]source}{}{}{}}
\makeatother
\crefname{originallemma}{original lemma}{original lemmas}
\Crefname{originallemma}{Original lemma}{Original lemmas}
\makeatletter
\newlabel{lemma-summands-relative-pseudo-coherent}{{09UN}{source}{}{sourcepointer-6}{}}
\newlabel{lemma-summands-relative-pseudo-coherent@cref}{{[originallemma][0][]09UN}{[1][0][]source}{}{}{}}
\makeatother
\crefname{originalreference}{original reference}{original references}
\Crefname{originalreference}{Original reference}{Original references}
\makeatletter
\newlabel{perfect-lemma-cohomology-base-change}{{08EU}{source}{}{sourcepointer-7}{}}
\newlabel{perfect-lemma-cohomology-base-change@cref}{{[originalreference][0][]08EU}{[1][0][]source}{}{}{}}
\makeatother
\title{Relative pseudo-coherence descends along an fppf covering}
\author{The Stacks Project Authors}
\date{stacks-project:more-morphisms.tex:local-snapshot-20261011; GNU FDL 1.2 or later}
\begin{document}
\maketitle
\noindent\textit{Editorial note.} Counted span is the author's own proof environment, source lines 20639-20680, immediately following the statement at 20630-20637.\par\medskip
\noindent\textit{Editorial note.} One cross-chapter pointer (Derived Categories of Schemes) plus six same-chapter lemmas are recorded in assembly.external\_references.\par\medskip
\noindent\textit{Editorial note.} The construction of the finite locally free pi used in the last step is displayed inside the span.\par\medskip
\noindent\textit{Allowed standard prerequisites (editorial).}\begin{itemize}
\item Flat morphisms of finite presentation are pseudo-coherent. chapter lemma lemma-flat-finite-presentation-pseudo-coherent, cited at proof lines 20641-20642. Makes each f\_i pseudo-coherent so that the pullback direction follows.
\item Dominating an fppf covering by a nicer one. chapter lemma lemma-dominate-fppf, cited in the proof. Replaces the covering by one whose pieces are better behaved, which is what the descent step needs.
\item Pulling back and summing relative pseudo-coherence. chapter lemmas lemma-pull-relative-pseudo-coherent and lemma-summands-relative-pseudo-coherent, cited in the proof. Move the condition along the covering maps and recognise a direct summand.
\item Finiteness and permanence inputs. chapter lemmas lemma-finite-morphism-relative-pseudo-coherence and lemma-permanence-pseudo-coherent with Derived Categories of Schemes, lemma perfect-lemma-cohomology-base-change, cited in the proof. Keep relative pseudo-coherence along the finite locally free morphism pi : Y -> X used at the end of the argument.
\end{itemize}\medskip
\ifdefined\hypertarget\hypertarget{sourcepointer-1}{}\fi
\noindent\textit{Source reference.} Original lemma 05WN: sources/stacks-more-morphisms/pinned-original-source.tex, source lines 13649--13649. Navigation pointer only: the counted proof cites this statement by \textbackslash{}ref as a standing input; the deduction the counted proof has to make is carried out inside the counted span, and no part of that deduction is delegated to this pointer.\par\medskip
\ifdefined\hypertarget\hypertarget{sourcepointer-2}{}\fi
\noindent\textit{Source reference.} Original lemma 09UK: sources/stacks-more-morphisms/pinned-original-source.tex, source lines 16995--16995. Navigation pointer only: the counted proof cites this statement by \textbackslash{}ref as a standing input; the deduction the counted proof has to make is carried out inside the counted span, and no part of that deduction is delegated to this pointer.\par\medskip
\ifdefined\hypertarget\hypertarget{sourcepointer-3}{}\fi
\noindent\textit{Source reference.} Original lemma 0695: sources/stacks-more-morphisms/pinned-original-source.tex, source lines 17428--17428. Navigation pointer only: the counted proof cites this statement by \textbackslash{}ref as a standing input; the deduction the counted proof has to make is carried out inside the counted span, and no part of that deduction is delegated to this pointer.\par\medskip
\ifdefined\hypertarget\hypertarget{sourcepointer-4}{}\fi
\noindent\textit{Source reference.} Original lemma 0683: sources/stacks-more-morphisms/pinned-original-source.tex, source lines 17440--17440. Navigation pointer only: the counted proof cites this statement by \textbackslash{}ref as a standing input; the deduction the counted proof has to make is carried out inside the counted span, and no part of that deduction is delegated to this pointer.\par\medskip
\ifdefined\hypertarget\hypertarget{sourcepointer-5}{}\fi
\noindent\textit{Source reference.} Original lemma 09US: sources/stacks-more-morphisms/pinned-original-source.tex, source lines 17155--17155. Navigation pointer only: the counted proof cites this statement by \textbackslash{}ref as a standing input; the deduction the counted proof has to make is carried out inside the counted span, and no part of that deduction is delegated to this pointer.\par\medskip
\ifdefined\hypertarget\hypertarget{sourcepointer-6}{}\fi
\noindent\textit{Source reference.} Original lemma 09UN: sources/stacks-more-morphisms/pinned-original-source.tex, source lines 17070--17070. Navigation pointer only: the counted proof cites this statement by \textbackslash{}ref as a standing input; the deduction the counted proof has to make is carried out inside the counted span, and no part of that deduction is delegated to this pointer.\par\medskip
\ifdefined\hypertarget\hypertarget{sourcepointer-7}{}\fi
\noindent\textit{Source reference.} Original reference 08EU: focused-seeds/geometry/S9-morph-rel-pseudo-coherent-descends-fppf/ref-tags.txt, source lines 5--5. Navigation pointer only: the counted proof cites this statement by \textbackslash{}ref as a standing input; the deduction the counted proof has to make is carried out inside the counted span, and no part of that deduction is delegated to this pointer.\par\medskip
\begin{lemma}
\label{lemma-relative-pseudo-coherent-descends-fppf}
Let $X \to S$ be locally of finite type.
Let $\{f_i : X_i \to X\}$ be an fppf covering of schemes.
Let $E \in D_\QCoh(\mathcal{O}_X)$. Let $m \in \mathbf{Z}$.
Then $E$ is $m$-pseudo-coherent relative to $S$
if and only if each $Lf_i^*E$ is $m$-pseudo-coherent relative to $S$.
\end{lemma}

\begin{proof}
Assume $E$ is $m$-pseudo-coherent relative to $S$.
The morphisms $f_i$ are pseudo-coherent by
Lemma \ref{lemma-flat-finite-presentation-pseudo-coherent}.
Hence $Lf_i^*E$ is $m$-pseudo-coherent relative to $S$
by Lemma \ref{lemma-pull-relative-pseudo-coherent}.

\medskip\noindent
Conversely, assume that $Lf_i^*E$ is $m$-pseudo-coherent relative to $S$
for each $i$. Pick $S = \bigcup U_j$, $W_j \to U_j$,
$W_j = \bigcup W_{j, k}$, $T_{j, k} \to W_{j, k}$, and
morphisms $\alpha_{j, k} : T_{j, k} \to X_{i(j, k)}$ over $S$ as in
Lemma \ref{lemma-dominate-fppf}.
Since the morphism $T_{j, K} \to S$ is flat and of finite presentation,
we see that $\alpha_{j, k}$ is pseudo-coherent by
Lemma \ref{lemma-permanence-pseudo-coherent}.
Hence
$$
L\alpha_{j, k}^*Lf_{i(j, k)}^*E = L(T_{i, k} \to S)^*E
$$
is $m$-pseudo-coherent relative to $S$ by
Lemma \ref{lemma-pull-relative-pseudo-coherent}.
Now we want to descend this property
through the coverings $\{T_{j, k} \to W_{j, k}\}$,
$W_j = \bigcup W_{j, k}$, $\{W_j \to U_j\}$, and $S = \bigcup U_j$.
Since for Zariski coverings the result is true
(by the definition of $m$-pseudo-coherence relative to $S$),
this means we may assume we have a single surjective
finite locally free morphism $\pi : Y \to X$
such that $L\pi^*E$ is pseudo-coherent relative to $S$.
In this case $R\pi_*L\pi^*E$ is pseudo-coherent relative to $S$
by Lemma \ref{lemma-finite-morphism-relative-pseudo-coherence}
(this is the first time we use that $E$ has quasi-coherent cohomology sheaves).
We have
$R\pi_*L\pi^*E = E \otimes^\mathbf{L}_{\mathcal{O}_X} \pi_*\mathcal{O}_Y$
for example by
Derived Categories of Schemes, Lemma
\ref{perfect-lemma-cohomology-base-change}
and locally on $X$ the map $\mathcal{O}_X \to \pi_*\mathcal{O}_Y$
is the inclusion of a direct summand. Hence we conclude by
Lemma \ref{lemma-summands-relative-pseudo-coherent}.
\end{proof}

\end{document}
